% SPDX-License-Identifier: CC-BY-4.0
% =====================================================================
% Chapter 8 of D. Joyner and Claude, "A hitch-hiker's guide to the
% metaplectic trace formula for SL(2): basic theory" (November 2026).
% New in Draft 3 (stage 10 of ../revision_log.md). It connects the text
% with the companion manuscript "Direct comparison of SL(2) with its
% metaplectic cover" (../../mult-one). It is new text, written by Claude,
% and has not been refereed. Checked by checks/comparison_n2_check.py.
% =====================================================================
\providecommand{\GA}{\overline{G(\A)}}
\providecommand{\Gb}{\overline{SL(2,F)}}

\chapter{Comparison with \texorpdfstring{$SL(2)$}{SL(2)}}\label{sec:compare}

The double cover lies over $SL(2)$. Its conjugacy classes, its tori and
its orbital integrals lie over those of $SL(2)$, and the two trace
formulas are indexed by the same rational classes. It is therefore
natural to lay the trace formula of
Chapters~\ref{sec:geometric}--\ref{sec:stable} beside the trace formula
of $SL(2)$ of Labesse and Langlands~\cite{LL79} and to compare them term
by term, an element $\gamma$ with a lift of the same $\gamma$. We call
this the \emph{direct comparison}. The companion text~\cite{JC26} studies
it from the side of local harmonic analysis. It shows that the direct
comparison works on the split torus and that it does not give a
correspondence of automorphic representations: the correspondence that
does exist, that of Shimura and Waldspurger, goes to $PGL(2)$ and
\emph{squares} the unramified parameters.

This chapter reads the same comparison through the formulas of the
earlier chapters. The two texts were written separately and neither used
the other. Put together, each explains something in the other. The
results of this text supply proofs of several statements that
\cite{JC26} quotes, and they locate the failure of the direct comparison
more precisely than \cite{JC26} does. In the other direction, the
question of~\cite{JC26} gives a use for the function $\Delta_\psi$ of
Lemma~\ref{lem:Delta} that Chapter~\ref{sec:stable} does not mention. Nothing in this
chapter is used in the earlier ones.

\paragraph{Results.}
\begin{itemize}
\item The matching of functions on the split torus used in~\cite{JC26} is
  the transfer of Lemma~\ref{lem:transfer} for Weil's character, and its
  factor is the restriction to the torus of the function $\Delta_\psi$
  that defines stable conjugacy (Lemma~\ref{lem:direct-torus}). So the
  direct comparison has a definition on all regular elements that does
  not depend on coordinates (Definition~\ref{def:direct}).
\item On the split torus the direct comparison is the character identity
  of Adams and Schultz, \eqref{eq:packet-lift}, in which the torus of
  $PGL(2)$ has been identified with the torus of $SL(2)$ by a map that
  does not come from a homomorphism of groups
  (Proposition~\ref{prop:wrong-orbit}). This is the square.
\item The function $\Delta_\psi$ satisfies a product formula on rational
  elements (Lemma~\ref{lem:Delta-product}). So if matching functions
  existed at every place, the regular elliptic and the regular hyperbolic
  orbital integrals of the two trace formulas would agree term by term
  (Corollary~\ref{cor:coherent}).
\item They do not exist. At an odd place, the stable orbital integral of
  the unit element $e_v$ of the Hecke algebra of the cover vanishes at
  the elliptic elements $\gamma$ of $K_v$ for which $2+\operatorname{tr}
  \gamma$ has odd valuation, and its $\kappa$-orbital integral does not.
  For the unit element of the Hecke algebra of $SL(2,F)$ it is the other
  way round (Proposition~\ref{prop:unit-swap}). So the fundamental lemma
  for the direct comparison fails on the elliptic set
  (Corollary~\ref{cor:no-fl}), and the rational elliptic terms of the two
  trace formulas differ by an explicit fourth root of unity
  (Proposition~\ref{prop:epsilon}).
\item On the spectral side, the group $GL(2,k)$ acts on the spectrum of
  the cover, and its orbits on the continuous and residual spectrum are
  the fibers of the squaring map (Proposition~\ref{prop:fibers}). The
  terms of the trace formula of $SL(2)$ that carry the multiplicities
  inside an $L$-packet have no counterpart on the cover
  (Section~\ref{sec:compare-spectral}).
\end{itemize}

\paragraph{Conventions.} A function on the linear group $SL(2,F)$ or
$G(\A)$ is written $\phi$, and a function on the cover is written $f$, as
in the rest of this text. (In~\cite{JC26} these are $f$ and
$\overline f$.) For $\phi\in C_c^\infty(SL(2,F))$ and a regular semisimple
$\gamma$ with centralizer $T_\gamma$ put
\[
J^\flat(\gamma,\phi)=|D(\gamma)|_v^{1/2}\int_{T_\gamma\backslash SL(2,F)}
\phi(y^{-1}\gamma y)\,dy,
\]
with the measures of Section~\ref{sec:kappa}. This is the orbital
integral $J_G$ of~\eqref{eq:JG} for the linear group. Its stable and
$\kappa$-versions are
\[
J^{\flat,\mathrm{st}}(\gamma,\phi)=\sum_{a\in F^\times/N E^\times}J^\flat(
\delta_a\gamma\delta_a^{-1},\phi),\qquad
J^{\flat,\kappa}(\gamma,\phi)=\sum_{a\in F^\times/N E^\times}\kappa_E(a)\,
J^\flat(\delta_a\gamma\delta_a^{-1},\phi),
\]
with $E=F[\gamma]$, as in~\cite{LL79}. Local statements are for $F=k_v$
and the fixed additive character $\psi=\psi_v$, as in
Chapter~\ref{sec:stable}.

\section{The two texts side by side}\label{sec:dictionary}

This section gives a dictionary between the notation of~\cite{JC26} and
the notation of this text, and lists the statements that \cite{JC26}
quotes and this text proves. The two texts use the same cocycle, Kubota's
$\beta$ of~\eqref{eq:beta}, so elements $(g,\zeta)$ of the local cover
mean the same thing in both.

\begin{center}
\renewcommand{\arraystretch}{1.15}
\begin{tabular}{p{0.44\textwidth}p{0.50\textwidth}}
In~\cite{JC26} & In this text\\\hline
$f$ on $SL(2)$, genuine $\overline f$ on the cover & $\phi$ on $SL(2)$,
  anti-genuine $f$ on the cover (for the double cover genuine and
  anti-genuine functions are the same)\\
genuine character $\overline\chi=\chi_\psi\mu$ of the cover of $A$, with
  $\mu$ a character of $F^\times$ & $\mu=\chi_\psi\cdot(\chi\circ p)$
  (Theorem~\ref{thm:corr}). The letters are exchanged: their $\mu$ is our
  $\chi$\\
unramified parameter $t'=\mu(\varpi)$ & $\chi_v(\varpi)$\\
the scalar $\overline c(s,\mu)$ of the intertwining operator &
  $m^S(\mu,\lambda)$ of~\eqref{eq:mS}, with $s=\lambda$\\
the splitting $\kappa_v(k)=(k,s_v(k))$ of $K_{0,v}$; ``original
  restricted-product coordinates'' & $K_v^*=\{(k,\kappa_v(k))\}$; the
  $\beta_v$-model of Section~\ref{sec:global}\\
odd matching on the hyperbolic set & the transfer of
  Lemma~\ref{lem:transfer} with $\bar\chi_0=\chi_\psi$\\
even matching, $h\mapsto(h,1)^2$ & the orbit map $t$
  of~\eqref{eq:packet-lift}, composed with $SL(2)\to PGL(2)$
  (Section~\ref{sec:wrong-orbit})\\
the theta series $\Theta_0$, and its complement ${\cal A}_{00}$ in the
  cusp forms & the cuspidal summands of the spaces $\Theta_b$ of
  Proposition~\ref{prop:theta-packet}(a), and the rest of
  $L^2_{\mathrm{cusp}}$\\
the Weil representation $r_{\psi,q}$ of~\cite{Fl80} and the function
  $\omega$ in its formula~(iii) & $\omega_{\psi'}$
  of~\eqref{eq:weil-rep}, and $\chi_{\psi'_{-1}}$
  (Lemma~\ref{lem:weil-rep})
\end{tabular}
\end{center}

\medskip
The text~\cite{JC26} marks the following statements as quoted from memory
or as conditional. They are proved here, for Kubota's cocycle and with
the constants determined.
\begin{enumerate}[label=(\roman*)]
\item \emph{Weil's product formula} and the fact that
  $\chi_\psi=\prod_v\chi_{\psi,v}$ is a genuine automorphic character of
  the cover of $A$: Lemma~\ref{lem:weil}(b),~(c). Weil's theorem itself
  is quoted from~\cite{We64}.
\item \emph{The global consistency of the direct comparison on the
  torus}: Lemma~\ref{lem:transfer} and the sentence after it. The two
  trace formulas of the torus and of its cover are the same Poisson
  summation formula.
\item \emph{The scalar of the intertwining operator},
  $L(2\lambda,\chi^2)/L(2\lambda+1,\chi^2)$:
  Lemma~\ref{lem:intertwining}(b), with the unramified computation done
  in the cover.
\item \emph{The residual spectrum}, one even Weil representation for each
  square class: Theorem~\ref{thm:residual}. Parts~(a) and~(b) and the
  matching of the parameters are proved here, and the irreducibility is
  quoted from~\cite{Ge76,GPS80,Wa80}.
\item \emph{The action of the torus in the Weil representation}, on which
  the exponent $\frac12$ in the proof that the theta series do not
  correspond to cusp forms of $SL(2)$ depends:
  Lemma~\ref{lem:weil-rep}. It maps $r^+_{\psi'}$ into
  $I(\chi_{\psi'_{-1}},-\frac12)$. The sign $-1$ in $\psi'_{-1}$ does
  not affect the absolute value, but it does affect which quadratic twist
  occurs, and it should be kept in mind when the normalization of
  Waldspurger's correspondence in~\cite{JC26} is checked.
\item \emph{The kernel and the character of the Weil representation},
  which \cite{JC26} computes formally and states as a conjecture:
  Lemma~\ref{lem:weil-kernel} gives the kernel with its constant, and
  Proposition~\ref{prop:weil-char} gives the character of
  $r^+_\psi-r^-_\psi$ on the regular set. The character of
  $r^+_\psi+r^-_\psi$ is~\eqref{eq:weil-char-full}. The device
  of~\cite{JC26} for passing between the two, the translation by $-1$, is
  the operation $f\mapsto f_-$ of Chapter~\ref{sec:stable}.
\end{enumerate}

\begin{example}\label{ex:dictionary}
Let $k=\Q$. In~\cite{JC26} the direct comparison is illustrated by the
function $\xi(2s)/\xi(1+2s)$, whose pole at $s=\frac12$ gives the theta
function $\sum_ne^{2\pi in^2z}$, against $\xi(s)/\xi(1+s)$ for $SL(2)$,
whose pole at $s=1$ gives the constants. Here these are
Example~\ref{ex:theta}: the residue of $m^S(\chi_\psi,\lambda)$ at
$\lambda=\frac12$ is $2/\pi^2$ for $S=\{\infty,2\}$, and the residual
representation is $V_1\cong r^+_{\psi_{-1}}$.
\end{example}

\section{The direct comparison and the function
\texorpdfstring{$\Delta_\psi$}{Delta\_psi}}\label{sec:direct}

This section shows that the factor in the matching of functions on the
split torus is a value of $\Delta_\psi$, and uses this to define the
direct comparison on all regular elements.

In~\cite{JC26} a function $\phi$ on $A(F)$ and a genuine function
$\overline\phi$ on the cover of $A(F)$ are said to correspond when their
Fourier transforms agree on corresponding characters. Globally the
choice that makes this consistent is $\overline\phi(a(x),1)=
\chi_{\psi}(a(x),1)^{-1}\phi(a(x))$. This is Lemma~\ref{lem:transfer} with
$\bar\chi_0=\chi_\psi$. The first observation is that the factor is not
special to the torus.

\begin{lemma}\label{lem:direct-torus}
Let $f\in C_c^\infty(\overline{A(F)})_-$, and let $\phi$ be the function
on $A(F)$ attached to $f$ by Lemma~\ref{lem:transfer} with
$\bar\chi_0=\chi_{\psi,v}$, so that $\phi(a)=f(\tilde a)\chi_{\psi,v}(
\tilde a)$ for any lift $\tilde a$ of $a$. Then
\[
f(\tilde a)=\Delta_\psi(\tilde a)\,\phi(a)\qquad\text{for every regular }
\tilde a\in\overline{A(F)}\text{ over }a .
\]
\end{lemma}

\begin{proof}
By Lemma~\ref{lem:Delta}(c), $\Delta_\psi$ is the complex conjugate of
$\chi_{\psi,v}$ on the regular elements of $\overline{A(F)}$. The values
of $\chi_{\psi,v}$ have absolute value $1$, so the conjugate is the
inverse.
\end{proof}

The function $\Delta_\psi$ is defined on all regular semisimple elements
of the cover, it is genuine, and it is invariant under conjugation
(Lemma~\ref{lem:Delta}(a)). This suggests the following definition.

\begin{definition}\label{def:direct}
Let $U$ be an open subset of the regular semisimple elements of
$SL(2,F)$ that is invariant under conjugation. Functions $\phi\in
C_c^\infty(SL(2,F))$ and $f\in C_c^\infty(\Gb)_-$ \emph{match directly on
$U$}, with respect to $\psi$, if
\[
J_G(x,f)=\Delta_\psi(x)\,J^\flat\bigl(p(x),\phi\bigr)\qquad\text{for all }x
\text{ with }p(x)\in U .
\]
\end{definition}

Both sides are genuine functions of $x$ that are invariant under
conjugation, so the condition makes sense. By
Lemma~\ref{lem:direct-torus}, on the hyperbolic set it is the odd
matching of~\cite{JC26}: the hyperbolic orbital integrals $J_G(\tilde a,f)$
and $J^\flat(a,\phi)$, as functions on the tori, correspond under
Lemma~\ref{lem:transfer}.

Why is a factor needed at all? One might try to compare $J_G$ at
$(\gamma,1)$ with $J^\flat$ at $\gamma$. But $\gamma\mapsto(\gamma,1)$ is
not compatible with conjugation: conjugating $(\gamma,1)$ by an element
of the cover gives $(\gamma',\pm1)$, with a sign that depends on the
elements. So the condition ``$J_G((\gamma,1),f)=J^\flat(\gamma,\phi)$ for all
$\gamma$'' is in general contradictory. The text~\cite{JC26} meets this
in its appendix on supercuspidal representations: the character relation
$\Theta_{\overline\pi}(t,\zeta)=\zeta\,\Theta_\pi(t)$ proposed in 1995
is not invariant under conjugation when $q\equiv3\pmod4$, and it is false
there. The function $\Delta_\psi$ is the correction. Its values are
fourth roots of unity, and $\Delta_\psi(\gamma,1)=1$ in the cases where
the uncorrected relation is plausible: for $F=\R$ and $\gamma$ a
rotation $k_\theta$, $0<\theta<\pi$ (Example~\ref{ex:Delta-real}), and
for $v$ odd, $\gamma\in K_v$ with lower left entry and
$2+\operatorname{tr}\gamma$ units (Example~\ref{ex:Delta-finite}).

\paragraph{Characters.} Let $\Theta$ be a genuine class function on
$\Gb$ and $\Theta^\flat$ a class function on $SL(2,F)$, both locally
integrable, with
\begin{equation}\label{eq:direct-char}
\Theta(x)=\overline{\Delta_\psi(x)}\;\Theta^\flat\bigl(p(x)\bigr)
\end{equation}
on the regular set. If $\phi$ and $f$ match directly on the whole
regular set, then $\int f\,\Theta=\int\phi\,\Theta^\flat$, by the Weyl
integration formula on the two groups. So~\eqref{eq:direct-char} is the
relation between characters that goes with
Definition~\ref{def:direct}. By Lemma~\ref{lem:Delta}(d) with $a=-1$,
and because $\Delta_\psi^2=(D^+,-1)_v$, the conjugate of $\Delta_\psi$
is $\Delta_{\psi_{-1}}$. By Proposition~\ref{prop:weil-char} the
relation~\eqref{eq:direct-char} therefore reads
\begin{equation}\label{eq:direct-weil}
\Theta(x)=|D^+(x)|_v^{1/2}\,\bigl(\Theta^+_{\psi_{-1}}(x)-\Theta^-_{\psi_{-1}}
(x)\bigr)\,\Theta^\flat\bigl(p(x)\bigr).
\end{equation}
In words: \emph{the direct comparison multiplies a character of $SL(2,F)$
by the stable character of the Weil representation, normalized by
$|\det(1+\gamma)|^{1/2}$.} This is an answer to the question
of~\cite{JC26} whether the odd matchings mean anything for
representations beyond the identities they were built to satisfy.

On the principal series the relation holds. Let
$\mu=\chi_{\psi,v}\cdot(\chi\circ p)$. By the proof of
Theorem~\ref{thm:IA}(a), the trace of a principal series is an integral
of hyperbolic orbital integrals against the inducing character,
$\operatorname{tr}I(\mu_\lambda,f)=\int_{\overline{A(F)}}\mu_\lambda(
\tilde a)J_G(\tilde a,f)\,d\tilde a$, and the same formula holds for the
principal series $I^\flat(\chi_\lambda)$ of $SL(2,F)$ induced from
$a(t)\mapsto\chi(t)|t|_v^\lambda$. If $\phi$ and $f$ match directly on the
hyperbolic set, the integrand is
\[
\chi_{\psi,v}(\tilde a)\,\chi_\lambda(a)\cdot\Delta_\psi(\tilde a)\,
J^\flat(a,\phi)=\chi_\lambda(a)\,J^\flat(a,\phi),
\]
because $\chi_{\psi,v}\Delta_\psi=1$ on the torus. It is invariant under
$\mu_2$, which has volume $1$. Hence
\begin{equation}\label{eq:direct-ps}
\operatorname{tr}I(\mu_\lambda,f)=\operatorname{tr}I^\flat(\chi_\lambda,\phi)
\qquad\text{for }\mu=\chi_{\psi,v}\cdot(\chi\circ p).
\end{equation}
In terms of characters: the character of $I(\mu_\lambda)$ at a regular
$\tilde a$ over $a$ is $(\mu_\lambda(\tilde a)+\mu^w_{-\lambda}(\tilde a))/
|D(a)|_v^{1/2}$, and $\mu^w=\chi_{\psi,v}\cdot(\chi^{-1}\circ p)$ by
Lemma~\ref{lem:w-action}(c). So it is $\chi_{\psi,v}(\tilde a)$ times the
character of $I^\flat(\chi_\lambda)$ at $a$, and both characters vanish on
the elliptic set. This is~\eqref{eq:direct-char} for
$\Theta=\Theta_{I(\mu_\lambda)}$ and $\Theta^\flat=\Theta_{I^\flat(
\chi_\lambda)}$.
The identity~\eqref{eq:direct-ps} is the identity of traces of principal
series of~\cite{JC26}.

\begin{example}\label{ex:direct-real}
Let $F=\R$ and $\psi(x)=e^{-2\pi ix}$. By Example~\ref{ex:weil-Q},
$\chi_\psi(a(t),1)$ is $1$ for $t>0$ and $i$ for $t<0$. So $\phi$ and $f$
match directly on the hyperbolic set when
$J_G((a(t),1),f)=J^\flat(a(t),\phi)$ for $t>0$ and $J_G((a(t),1),f)=-i\,
J^\flat(a(t),\phi)$ for $t<0$. This is the real matching
of~\cite{JC26}, in which the factor $-i$ appears for $x<0$. So the
choice made there is the one that belongs to the additive character
$e^{-2\pi ix}$ at the real place. For $e^{2\pi ix}$ the factor would be
$i$.
\end{example}

\section{The same factor, a different orbit}\label{sec:wrong-orbit}

This section compares the direct comparison with the character identity
that is known to hold between the cover and $SO(3)=PGL(2)$, and finds
that they differ only in the way the tori are identified.

The identity~\eqref{eq:packet-lift} of Adams, Schultz and Luo has the
same first factor as~\eqref{eq:direct-weil}, the stable character of the
Weil representation. In~\eqref{eq:packet-lift} the second factor is a
character of $PGL(2,F)$ evaluated at $t(x)$, the class with eigenvalues
$\epsilon,1,\epsilon^{-1}$ in $SO(3)$ when $p(x)$ has eigenvalues
$\epsilon^{\pm1}$. In~\eqref{eq:direct-weil} it is a character of
$SL(2,F)$ evaluated at $p(x)$ itself. The next proposition makes the
comparison exact on the split torus.

For a character $\chi$ of $F^\times$ let $\sigma(\chi)$ be the principal
series of $PGL(2,F)$ induced from $\diag(t_1,t_2)\mapsto\chi(t_1/t_2)$,
and $I^\flat(\chi)$ the principal series of $SL(2,F)$ induced from
$a(t)\mapsto\chi(t)$, both normalized.

\begin{proposition}\label{prop:wrong-orbit}
Let $s\in F^\times$, $s\ne\pm1$, and let $x$ be a lift of $a(s)$.
\begin{enumerate}[label=(\alph*)]
\item The class $t(x)$ is that of $\diag(s,1)$ in $PGL(2,F)$, and
  \[
  D\bigl(a(s)\bigr)=D^+\bigl(a(s)\bigr)\cdot D'(s),\qquad
  D'(s)=\frac{(1-s)^2}{s},
  \]
  where $D'(s)$ is the Weyl discriminant of $\diag(s,1)$ in $PGL(2,F)$.
\item $|D^+(x)|_v^{-1/2}\,\Theta_{\sigma(\chi)}\bigl(\diag(s,1)\bigr)=
  \Theta_{I^\flat(\chi)}\bigl(a(s)\bigr)$.
\item The restriction of $\sigma(\chi)$ to $SL(2,F)$ is $I^\flat(\chi^2)$.
\end{enumerate}
So on the split torus the right side of~\eqref{eq:packet-lift} for
$\sigma_v=\sigma(\chi)$ and $\psi'=\psi_{-1}$ is the right side
of~\eqref{eq:direct-weil} for $\Theta^\flat=\Theta_{I^\flat(\chi)}$, and both
are the character of $I(\chi_{\psi,v}\cdot(\chi\circ p))$.
\end{proposition}

\begin{proof}
(a) The eigenvalues of $a(s)$ are $s^{\pm1}$, and $\diag(s,1)$ acts on
the Lie algebra of $PGL(2)$ with eigenvalues $s,1,s^{-1}$. Next,
$D(a(s))=(s+s^{-1})^2-4=(s-s^{-1})^2=(s-1)^2(s+1)^2/s^2$ and
$D^+(a(s))=2+s+s^{-1}=(1+s)^2/s$. The Weyl discriminant of $\diag(s,1)$ is
$(1-s)(1-s^{-1})$ up to sign, that is, $(1-s)^2/s$.

(b) The character of a normalized principal series on a regular element
of the split torus is the sum of the inducing character over the Weyl
group, divided by the square root of the absolute value of the Weyl
discriminant. For $\sigma(\chi)$ at $\diag(s,1)$ this is
$(\chi(s)+\chi(s)^{-1})/|D'(s)|_v^{1/2}$, and for $I^\flat(\chi)$ at $a(s)$
it is $(\chi(s)+\chi(s)^{-1})/|D(a(s))|_v^{1/2}$. Apply~(a).

(c) The torus of $SL(2)$ maps to the torus of $PGL(2)$ by
$a(t)\mapsto\diag(t,t^{-1})=\diag(t^2,1)$ modulo the center, and the
inducing character $\diag(t_1,t_2)\mapsto\chi(t_1/t_2)$ takes the value
$\chi(t^2)$ there. Restriction from $PGL(2,F)$ to the image of $SL(2,F)$
commutes with induction from the Borel subgroups, because
$PGL(2,F)=B'(F)\cdot\operatorname{im}SL(2,F)$ for the Borel subgroup
$B'$ of $PGL(2)$.

The last statement follows from~(b) and the computation of the
character of $I(\mu_\lambda)$ in Section~\ref{sec:direct}.
\end{proof}

The proposition says where the square comes from. The correct
comparison, \eqref{eq:packet-lift}, pairs the genuine principal series
with parameter $\chi$ with the representation $\sigma(\chi)$ of
$PGL(2,F)$. The elements that are compared are a lift of $a(s)$ and
$\diag(s,1)$. To get from $PGL(2)$ to $SL(2)$ there are two things one
can do, and they differ.
\begin{itemize}
\item One can \emph{restrict}. This is what the automorphic
  correspondence does. By~(c) the result is $I^\flat(\chi^2)$, and the
  element $a(s)$ of $SL(2)$ goes to $\diag(s^2,1)$, which is compared
  with a lift of $a(s^2)$. This is the squaring correspondence, or even
  matching, of~\cite{JC26}.
\item One can \emph{identify the tori}, sending $\diag(s,1)$ to $a(s)$.
  By~(b) the result is $I^\flat(\chi)$, and $a(s)$ is compared with a lift
  of $a(s)$. This is the direct comparison. The map
  $\diag(s,1)\mapsto a(s)$ is an isomorphism of tori that commutes with
  the Weyl groups, but it is not the restriction of a homomorphism
  between $PGL(2)$ and $SL(2)$ in either direction.
\end{itemize}
So the direct comparison uses the right transfer factor and the wrong
correspondence of orbits. On the dual side this is the statement
of~\cite{JC26} that the direct comparison identifies the dual tori of
$SL(2,\C)$ and $PGL(2,\C)$ by $t=\pm t'$, while the only homomorphism
between the dual groups gives $t=t'^{\,2}$ (compare the paragraph on the
dual group in Section~\ref{sec:stable-conj}).

The squares in a torus have a special position in
Chapter~\ref{sec:stable} as well. If $\gamma=\delta^2$ with $\delta$ in
the same torus, then $1+\gamma=\delta(\delta+\delta^{-1})$ and
$\delta+\delta^{-1}$ is the scalar $\operatorname{tr}\delta$, so
\begin{equation}\label{eq:square-Dplus}
D^+(\delta^2)=\det(1+\delta^2)=(\operatorname{tr}\delta)^2 .
\end{equation}
Hence $\eta_{\delta^2}=1$: on the elements of the cover that lie over
squares, the twisted action $\star$ of~\eqref{eq:star} is ordinary
conjugation by $GL(2,F)$, and conjugation by $GL(2,F)$ never exchanges
the two lifts. Away from them it does. By
Lemma~\ref{lem:kappa-orbital}(b) the average of $J_G(x,f^h)$ over
$h\in GL(2,F)$ vanishes unless $D^+(\gamma)$ is a square or $\theta$
times a square. A character of $PGL(2,F)$ restricted to $SL(2,F)$ is
invariant under $GL(2,F)$. So a comparison of the cover with $SL(2)$
that passes through $PGL(2)$ can only see the cover near the elements
over squares, and this is the support condition in the definition of an
even matching in~\cite{JC26}.

\begin{example}\label{ex:wrong-orbit}
Let $v$ be odd, $\chi$ unramified and $\chi(\varpi)=t'$. The genuine
unramified principal series with this parameter corresponds, by the
Shimura--Waldspurger correspondence, to the representation $\sigma(\chi)$
of $PGL(2,F)$ with Satake parameter $\diag(t',t'^{-1})\in SL(2,\C)$. Its
restriction to $SL(2,F)$ is $I^\flat(\chi^2)$, with parameter $t=t'^{\,2}$.
The direct comparison assigns $I^\flat(\chi)$, with $t=t'$, or
$I^\flat(\chi\kappa_u)$ with $t=-t'$ if $\psi$ is replaced by $\psi_u$, $u$
a unit that is not a square (Lemma~\ref{lem:Delta}(d)). For
$\Delta$, the cusp form of weight $12$, with $\tau(p)=p^{11/2}(\alpha_p+
\alpha_p^{-1})$, the form of weight $13/2$ attached to it has
$t'=\pm\alpha_p$, and the restriction of $\Delta$ to $SL(2)$ has
$t=\alpha_p^2$.
\end{example}

\section{Rational elements}\label{sec:compare-rational}

This section proves that the factor $\Delta_\psi$ is coherent on
rational elements, and draws the consequence for the geometric sides.

Let $\gamma\in G(\Q)=SL(2,k)$ be regular semisimple. Its image
$s(\gamma)$ in $\GA$ has local components, elements $x_v$ of
$\overline{SL(2,k_v)}$ over $\gamma$ whose product is $s(\gamma)$. They
are determined up to signs $\zeta_v\in\mu_2$, almost all $1$, with
$\prod_v\zeta_v=1$. Since each $\Delta_{\psi_v}$ is genuine, the product
$\prod_v\Delta_{\psi_v}(x_v)$ does not depend on the choice, provided
that almost all factors are $1$.

\begin{lemma}\label{lem:Delta-product}
Let $\gamma\in SL(2,k)$ be regular semisimple, and let $(x_v)$ be local
components of $s(\gamma)$. Then $\Delta_{\psi_v}(x_v)=1$ for almost all
$v$, and
\[
\prod_v\Delta_{\psi_v}(x_v)=1 .
\]
\end{lemma}

\begin{proof}
Both statements are unchanged when $\gamma$ is replaced by a conjugate
under $SL(2,k)$, because $s$ is a homomorphism and each $\Delta_{\psi_v}$
is invariant under conjugation. So we may assume that the lower left
entry $c$ of $\gamma$ is not $0$. By Lemma~\ref{lem:global}(a),
$s(\gamma)=(\gamma,\kappa_\A(\gamma))$ in the coordinates of the cocycle
$\sigma_\A$, with $\kappa_\A=\prod_v\kappa_v$. So we may take $x_v=(\gamma,
\kappa_v(\gamma))$ in the $\sigma_v$-model. Under the identification
$(g,\zeta)_\sigma\mapsto(g,\kappa_v(g)\zeta)_\beta$ of
Section~\ref{sec:global} this is the element $(\gamma,1)$ of the
$\beta_v$-model, in which~\eqref{eq:Delta} is written. Hence
\[
\Delta_{\psi_v}(x_v)=\bigl(c,D^+(\gamma)\bigr)_v\,\frac{\gamma_v(D^+(\gamma))}
{\gamma_v(1)} .
\]
Here $c$ and $D^+(\gamma)=2+\operatorname{tr}\gamma$ are elements of
$k^\times$. At almost all $v$ both are units, $v$ is non-dyadic and
$\psi_v$ is unramified, so the Hilbert symbol is $1$, and the Weil
indices are $1$ by Lemma~\ref{lem:weil}(b). The product of the Hilbert
symbols over all $v$ is $1$ by~\ref{H4}, and the products of
$\gamma_v(D^+(\gamma))$ and of $\gamma_v(1)$ are $1$ by
Lemma~\ref{lem:weil}(b).
\end{proof}

For the split torus the lemma is the statement that $\chi_\psi$ is trivial
on $A(\Q)$ (Lemma~\ref{lem:weil}(c)), by Lemma~\ref{lem:Delta}(c). In the
general case it has the same meaning: $\Delta_\psi=\prod_v\Delta_{\psi_v}$
is a genuine function on the regular elements of $\GA$ that is $1$ on
$s(G(\Q))$, so multiplying by it carries functions on $G(\A)$ to genuine
functions on $\GA$ without disturbing the rational points.

\begin{corollary}\label{cor:coherent}
Let $\phi=\bigotimes_v\phi_v$ on $G(\A)$ and $f=\bigotimes_vf_v\in
\mathcal H(\GA)$ be such that $\phi_v$ and $f_v$ match directly on the
regular set at every place $v$, with respect to $\psi_v$. Then for every
regular semisimple $\gamma\in G(\Q)$
\[
J_G\bigl(s(\gamma),f\bigr)=J^\flat(\gamma,\phi),
\]
the products of the local orbital integrals. Consequently the regular
elliptic terms of the two trace formulas agree class by class, and so do
the unweighted orbital integrals at the regular hyperbolic classes.
\end{corollary}

\begin{proof}
$J_G(s(\gamma),f)=\prod_vJ_G(x_v,f_v)=\prod_v\Delta_{\psi_v}(x_v)J^\flat(\gamma,
\phi_v)$, and the product of the factors $\Delta_{\psi_v}(x_v)$ is $1$ by
Lemma~\ref{lem:Delta-product}. By Proposition~\ref{prop:elliptic} the
term of a regular elliptic class is $\vol(T_\gamma(\Q)\backslash
T_\gamma(\A))$ times this orbital integral, and the same holds for
$SL(2)$ with the same volume.
\end{proof}

The corollary is a statement about functions that match at \emph{every}
place. In a comparison of trace formulas one has no choice at almost all
places: $\phi_v$ and $f_v$ must be the unit elements of the Hecke
algebras, or at least correspond under an isomorphism of Hecke algebras
that is compatible with the comparison of the unramified
representations. On the hyperbolic set the unit elements do match
directly. The next section shows that on the elliptic set they do not.

\begin{example}\label{ex:coherent}
Let $k=\Q$ and
\[
\gamma=\begin{pmatrix}2&-1\\3&-1\end{pmatrix}\in SL(2,\Z),\qquad
\operatorname{tr}\gamma=1,\quad D^+(\gamma)=3,\quad D(\gamma)=-3 .
\]
It is elliptic, of order $6$, with $k[\gamma]=\Q(\sqrt{-3})$. In the
$\beta_v$-model, with $\psi$ as in Example~\ref{ex:weil-Q},
\[
\Delta_{\psi_\infty}(\gamma,1)=1,\qquad\Delta_{\psi_2}(\gamma,1)=i,\qquad
\Delta_{\psi_3}(\gamma,1)=-i,
\]
and $\Delta_{\psi_p}(\gamma,1)=1$ for $p\ge5$, because $c=3$ and $D^+=3$
are units there. The product is $1$. At $p=3$ we have $(3,3)_3=(3,-1)_3=-1$
and $\gamma_3(3)=i$ (Example~\ref{ex:weil-Q}). The script computes all
three values.
\end{example}

\section{The unit elements on the elliptic set}\label{sec:unit-elliptic}

This section computes the stable and $\kappa$-orbital integrals of the
unit elements of the two Hecke algebras at the elliptic elements of
$K_v$ that are close to $-1$, and shows that they are interchanged.

Let $v$ be finite and non-dyadic, $F=k_v$, $\psi=\psi_v$ unramified,
$K=K_v$, and let $e_v$ be the unit element~\eqref{eq:unit} of the Hecke
algebra of the cover, supported on $\overline K$ and equal to $1$ on
$K^*$. Let $\mathbf 1_K$ be the characteristic function of $K$ on
$SL(2,F)$. For $\gamma\in K$ let $\gamma^*=(\gamma,\kappa_v(\gamma))$ be
its lift in $K^*$ ($\beta_v$-model).

\begin{proposition}\label{prop:unit-swap}
Let $\gamma\in K$ be regular semisimple, and let $u\in\OO_v^\times$ be a
unit that is not a square.
\begin{enumerate}[label=(\alph*)]
\item If $D^+(\gamma)=2+\operatorname{tr}\gamma$ is a unit, then
  $\Delta_\psi(\gamma^*)=1$.
\item If $|D(\gamma)|_v=1$, or if $D(\gamma)$ has valuation $1$, then
  \[
  J_G(\gamma^*,e_v)=J^\flat(\gamma,\mathbf 1_K)=|D(\gamma)|_v^{1/2}\,
  \vol(T_\gamma\cap K)^{-1}>0 .
  \]
\item Suppose that $D^+(\gamma)$ has odd valuation. Then $\gamma$ is
  elliptic, $E=F[\gamma]$ is ramified, $F^\times/N E^\times=\{1,u\}$,
  and $\eta_\gamma(u)=\kappa_E(u)=-1$. Moreover
  \[
  J_G(\delta_u\star\gamma^*,e_v)=-J_G(\gamma^*,e_v),\qquad
  J^\flat(\delta_u\gamma\delta_u^{-1},\mathbf 1_K)=J^\flat(\gamma,\mathbf 1_K).
  \]
  Hence
  \begin{align*}
  J^{\mathrm{st}}_G(\gamma^*,e_v)&=0,&
  J^{\kappa}_G(\gamma^*,e_v)&=2J_G(\gamma^*,e_v),\\
  J^{\flat,\mathrm{st}}(\gamma,\mathbf 1_K)&=2J^\flat(\gamma,\mathbf 1_K),&
  J^{\flat,\kappa}(\gamma,\mathbf 1_K)&=0 .
  \end{align*}
\item In the situation of~(c), $\Delta_\psi(\gamma^*)^2=(-1,\varpi)_v$.
  So $\Delta_\psi(\gamma^*)=\pm i$ if $q_v\equiv3\pmod4$ and
  $\Delta_\psi(\gamma^*)=\pm1$ if $q_v\equiv1\pmod4$, and the sign changes
  when $\gamma$ is replaced by $\delta_u\gamma\delta_u^{-1}$.
\end{enumerate}
\end{proposition}

\begin{proof}
(a) Write $\gamma=\begin{pmatrix}\alpha&b\\c&d\end{pmatrix}$. Since $K^*$ is
a subgroup and $\Delta_\psi$ is invariant under conjugation, we may
replace $\gamma$ by a conjugate under $K$. If $c=b=0$, then
$\gamma=a(t)$ with $t$ a unit, $\gamma^*=(a(t),1)$, and
$\Delta_\psi(\gamma^*)=\gamma_v(1)/\gamma_v(t)=1$ by
Lemmas~\ref{lem:Delta}(c) and~\ref{lem:weil}(b). If $c=0\ne b$, conjugate
by $w$. So let $c\ne0$. The Weil indices in~\eqref{eq:Delta} are $1$,
because $D^+$ is a unit. If $c$ is a unit, then $\kappa_v(\gamma)=1$ and
$(c,D^+)_v=1$. If $0<|c|_v<1$, then $\kappa_v(\gamma)=(c,d)_v$ and
$\Delta_\psi(\gamma^*)=(c,d\,D^+)_v$. Now $dD^+=d^2+2d+\alpha d=(d+1)^2+bc$,
so $dD^+\equiv(d+1)^2$ modulo $\varpi$. Since $dD^+$ is a unit, so is
$d+1$, and $dD^+$ is a square in $\OO_v^\times$. Hence the symbol is $1$.

(b) If $y^{-1}\gamma y\in K$ with $y\in SL(2,F)$, then $\gamma$ stabilizes
the lattice $y\OO_v^2$. The ring $\OO_v[\gamma]$ has discriminant
$D(\gamma)$. Under either hypothesis this is the discriminant of the
maximal order of the algebra $E=F[\gamma]$, so $\OO_v[\gamma]=\OO_E$. The
lattices in $F^2$ that are stable under $\OO_E$ are the $e\,\OO_v^2$ with
$e\in E^\times$, because $F^2$ is free of rank one over $E$ and $\OO_E$ is
a principal ideal ring (a product of two such, if $E=F\times F$). So
$y=ek$ with $e\in E^\times$ and $k\in GL(2,\OO_v)$, and $\det e$ is a
unit. The elements of $E^\times$ of unit norm are, modulo $\OO_E^\times$,
the powers of an element of $T_\gamma$ in the split case and trivial in
the other cases; so $y\in T_\gamma\cdot K$. Therefore the orbital
integrals are integrals over $(T_\gamma\cap K)\backslash K$, where
$\mathbf 1_K(k^{-1}\gamma k)=1$ and $e_v(\tilde k^{-1}\gamma^*\tilde k)=1$
because $K^*$ is a group. With $\vol(K)=1$ this gives the value stated.

(c) Since $D^+$ is not a unit and $4$ is, $D^+-4$ is a unit, and
$D(\gamma)=(\operatorname{tr}\gamma)^2-4=D^+(D^+-4)$ has the same
valuation as $D^+$, which is odd. So $E=F(\sqrt{D(\gamma)})$ is a ramified
quadratic extension. The units that are norms from a ramified extension
are the squares, so $u\notin N E^\times$ and $\kappa_E(u)=-1$. Also
$\eta_\gamma(u)=(u,D^+)_v=-1$, because $u$ is a non-square unit and $D^+$
has odd valuation, by~\ref{H3}. By~\eqref{eq:star},
$\delta_u\star\gamma^*=-\Ad(\delta_u)\gamma^*$. Now $J_G(\Ad(\delta_u)
\gamma^*,e_v)=J_G(\gamma^*,e_v^{\delta_u})$ and $e_v^{\delta_u}=e_v$ by
Lemma~\ref{lem:Ad}(e). Since $e_v$ is anti-genuine,
$J_G(\delta_u\star\gamma^*,e_v)=-J_G(\gamma^*,e_v)$. On $SL(2,F)$,
conjugation by $\delta_u$ preserves $K$ and the measures, which gives the
second equality. The four values follow from~\eqref{eq:Jst} and the
definitions of $J^{\flat,\mathrm{st}}$ and $J^{\flat,\kappa}$, with
$a\in\{1,u\}$.

(d) By~\eqref{eq:Delta} and Lemma~\ref{lem:weil}(a),
$\Delta_\psi^2=(D^+,-1)_v$, which is $(-1,\varpi)_v$ when $D^+$ has odd
valuation. By Lemma~\ref{lem:Ad}(e), $\Ad(\delta_u)\gamma^*$ is the lift
in $K^*$ of $\delta_u\gamma\delta_u^{-1}$, and $\Delta_\psi(\Ad(\delta_u)
\gamma^*)=\eta_\gamma(u)\Delta_\psi(\gamma^*)$ by Lemma~\ref{lem:Delta}(b).
\end{proof}

Part~(c) is the main point. For the linear group, the two conjugacy
classes in the stable class of $\gamma$ both meet $K$, and $\mathbf 1_K$
has the same orbital integral on them. So its stable orbital integral is
twice that number and its $\kappa$-orbital integral is $0$. For the
cover the two classes in the stable class of $\gamma^*$ are those of
$\gamma^*$ and of $\delta_u\star\gamma^*$. The second is \emph{minus} the
lift in $K^*$ of $\delta_u\gamma\delta_u^{-1}$: the twist that makes
stable conjugacy coherent at the rational points moves the element out
of $K^*$. So the unit element takes opposite values on the two classes.
Its stable orbital integral is $0$ and its $\kappa$-orbital integral is
twice the orbital integral.

This agrees with Lemma~\ref{lem:kappa-orbital}(c). The $\kappa$-orbital
integral at $\gamma^*$ is the stable orbital integral at
$\tilde z\gamma^*$, which lies over $-\gamma$, and $D^+(-\gamma)=
4-D^+(\gamma)$ is a unit: at $-\gamma$ nothing unusual happens. The
elements in question are those congruent to $-1$ modulo the prime of
$E$. Near $+1$ the cover is a product and the direct comparison is the
identity. Near $-1$ the stable and the $\kappa$-orbital integrals of the
unit element are interchanged, relative to $SL(2)$.

\begin{remarkrev}\label{rem:li-fl}
\emph{Comparison with the fundamental lemma for $SO(3)$.} The vanishing
of $J^{\mathrm{st}}_G(\gamma^*,e_v)$ is what the fundamental lemma
of~\cite{Li11} predicts. Write $\gamma=e/\bar e$ with $e\in E^\times$
(Hilbert~90). The class in $PGL(2,F)=SO(3)$ that corresponds to
$\gamma^*$ is the image of $e$, whose eigenvalues in $SO(3)$ are
$\gamma,1,\gamma^{-1}$. Now $1+\gamma=\operatorname{tr}(e)/\bar e$, so
$D^+(\gamma)=\operatorname{tr}(e)^2/N(e)$, and $D^+(\gamma)$ has odd
valuation exactly when $\det e=N(e)$ does. An element of $GL(2,F)$ whose
determinant has odd valuation is not conjugate into the center times
$GL(2,\OO_v)$, so the unit element of the Hecke algebra of $PGL(2,F)$
has orbital integral $0$ at its image. We have not compared the
normalizations of the transfer factor of~\cite{Li11} with ours, so this
is a consistency check and not a second proof.
\end{remarkrev}

\begin{corollary}\label{cor:no-fl}
Let $\gamma\in K$ be regular with $D^+(\gamma)$ of odd valuation. There
is no genuine function $\Delta'$ on the elements of $\Gb$ over the stable
class of $\gamma$ that is nowhere $0$, invariant under the action $\star$
of $GL(2,F)$, and satisfies
\[
J_G(x,e_v)=\Delta'(x)\,J^\flat\bigl(p(x),\mathbf 1_K\bigr)
\]
there. In particular, since $\Delta_\psi$ is such a function, the unit
elements $\mathbf 1_K$ and $e_v$ do not match directly, in the sense of
Definition~\ref{def:direct}, on any invariant open set that contains
$\gamma$.
\end{corollary}

\begin{proof}
The two conjugacy classes in the stable class of $\gamma^*$ are those of
$\gamma^*$ and $\delta_u\star\gamma^*$, and $\Delta'$ takes the same value
on them. Summing the identity over the two classes gives
$J^{\mathrm{st}}_G(\gamma^*,e_v)=\Delta'(\gamma^*)\,J^{\flat,\mathrm{st}}(
\gamma,\mathbf 1_K)$. By Proposition~\ref{prop:unit-swap}(c) the left side
is $0$ and $J^{\flat,\mathrm{st}}(\gamma,\mathbf 1_K)=2J^\flat(\gamma,
\mathbf 1_K)$, which is positive because $\gamma\in K$. So
$\Delta'(\gamma^*)=0$, a contradiction. The function $\Delta_\psi$ is
genuine, has values in $\mu_4$, and is invariant under $\star$ by
Lemma~\ref{lem:Delta}(b).
\end{proof}

So the direct comparison has no fundamental lemma on the elliptic set.
On the hyperbolic set it has one: the unit elements correspond under the
Satake isomorphism of~\cite{JC26}, and Proposition~\ref{prop:unit-swap}(a)
shows $\Delta_\psi=1$ on the elements $(a(t),1)$, $t$ a unit, where
the Satake transforms of both unit elements are supported. The
text~\cite{JC26} found that the comparison is consistent on the split
torus and fails for the terms of the spectral side that involve the
root, and it left the elliptic terms unexamined. The proposition settles
that point: the elliptic terms fail too, and they fail at the elements
congruent to $-1$.

The failure can be measured on the rational elements.

\begin{proposition}\label{prop:epsilon}
Let $S$ be a finite set of places of $\Q$ containing $\infty$ and $2$,
such that $\psi_v$ is unramified for every place $v$ of $k$ outside
$S_k$. Let $\phi=\phi_S\otimes\bigotimes_{v\notin S_k}\mathbf 1_{K_v}$ and
$f=f_S\otimes e^S$, where $\phi_v$ and $f_v$ match directly on the regular
set for each $v\in S_k$. Let $\gamma\in SL(2,k)$ be regular semisimple,
with $\gamma\in K_v$ and $D(\gamma)$ of valuation at most $1$ for every
$v\notin S_k$. Then
\[
J_G\bigl(s(\gamma),f\bigr)=\varepsilon^S(\gamma)\,J^\flat(\gamma,\phi),\qquad
\varepsilon^S(\gamma)=\prod_{v\notin S_k}\Delta_{\psi_v}(\gamma^*_v)^{-1},
\]
where $\gamma_v^*$ is the lift of $\gamma$ in $K_v^*$. The factor
$\varepsilon^S(\gamma)$ is a fourth root of unity. It is $1$ if
$D^+(\gamma)$ is a unit outside $S_k$. It is not $1$ in general.
\end{proposition}

\begin{proof}
Take the local components $x_v=(\gamma,1)$ of $s(\gamma)$ in the
$\beta_v$-models, as in the proof of Lemma~\ref{lem:Delta-product}. For
$v\in S_k$, $J_G(x_v,f_v)=\Delta_{\psi_v}(x_v)J^\flat(\gamma,\phi_v)$. For
$v\notin S_k$, $x_v=\kappa_v(\gamma)\gamma_v^*$, so
$J_G(x_v,e_v)=\kappa_v(\gamma)J_G(\gamma^*_v,e_v)=\kappa_v(\gamma)J^\flat(
\gamma,\mathbf 1_{K_v})$ by Proposition~\ref{prop:unit-swap}(b). Multiply,
and use Lemma~\ref{lem:Delta-product} in the form $\prod_{v\in S_k}
\Delta_{\psi_v}(x_v)=\prod_{v\notin S_k}\Delta_{\psi_v}(x_v)^{-1}$. Since
$\Delta_{\psi_v}(x_v)=\kappa_v(\gamma)\Delta_{\psi_v}(\gamma^*_v)$ and
$\kappa_v(\gamma)=\pm1$, the factor at $v\notin S_k$ is
$\Delta_{\psi_v}(\gamma_v^*)^{-1}$. It is $1$ where $D^+(\gamma)$ is a unit,
by Proposition~\ref{prop:unit-swap}(a). An example with
$\varepsilon^S(\gamma)\ne1$ follows.
\end{proof}

For the hyperbolic classes no restriction on $D(\gamma)$ is needed. If
$\gamma=a(t)$ with $t$ a unit outside $S_k$, then at each $v\notin S_k$
the hyperbolic orbital integrals of $\mathbf 1_{K_v}$ and of $e_v$ are
their Satake transforms, which are the characteristic function of
$A(\OO_v)$ and its genuine extension from $A(\OO_v)\times\{1\}$. Since
$\Delta_{\psi_v}=1$ on $A(\OO_v)\times\{1\}$, the unit elements match
directly on the hyperbolic set, and $J_G(s(\gamma),f)=J^\flat(\gamma,\phi)$
by the proof of Corollary~\ref{cor:coherent}. This is the consistency of
the direct comparison on the split torus.

\begin{example}\label{ex:epsilon}
Let $k=\Q$, $S=\{\infty,2\}$ and $\gamma=\begin{pmatrix}2&-1\\3&-1
\end{pmatrix}$, as in Example~\ref{ex:coherent}. Here $D(\gamma)=-3$, so
the hypothesis holds, and only $p=3$ contributes. We have
$\kappa_3(\gamma)=(3,-1)_3=-1$, so
$\Delta_{\psi_3}(\gamma_3^*)=-\Delta_{\psi_3}(\gamma,1)=i$ and
\[
\varepsilon^S(\gamma)=-i .
\]
For any test functions that match directly at $\infty$ and at $2$, the
term of $\gamma$ in the trace formula of the cover is $-i$ times its
term in the trace formula of $SL(2)$. The element
$\gamma'=\delta_{-1}\gamma\delta_{-1}^{-1}$ is conjugate to $\gamma$
under $GL(2,\Q)$ and not under $SL(2,\Q)$, since $-1$ is not a norm from
$\Q_3(\sqrt{-3})$. For it $\varepsilon^S(\gamma')=+i$, by
Proposition~\ref{prop:unit-swap}(d). So the two classes contribute
$-i\,J^\flat(\gamma,\phi)$ and $+i\,J^\flat(\gamma',\phi)$ to the trace
formula of the cover. Up to the factor $-i$, their sum is the
\emph{difference} of the two orbital integrals of $\phi$, where the trace
formula of $SL(2)$ has the sum. The sign
$\kappa_3(\gamma)=-1$ is a value of the Jacobi symbol of
Example~\ref{ex:jacobi}, the multiplier of the theta function. This is
how the multiplier system of half-integral weight enters the comparison.
\end{example}

\section{The spectral side}\label{sec:compare-spectral}

This section goes through the spectral side of
Corollary~\ref{cor:itf-final} and says, for each kind of term, what the
direct comparison does with it and what the squaring correspondence
does. The statements about the trace formula of $SL(2)$ are quoted
from~\cite{LL79} and~\cite{GJ79} and were not rechecked
(Remark~\ref{rem:compare-status}).

\paragraph{The continuous spectrum.} By~\eqref{eq:direct-ps}, directly
matching functions satisfy $\operatorname{tr}I(\mu_{it},f)=
\operatorname{tr}I^\flat(\chi_{it},\phi)$ for $\mu=\chi_\psi\cdot(\chi\circ
p)$. In the trace formula these traces are integrated against the
logarithmic derivative of the normalizing factor. For the cover this is
$m(\mu,\lambda)$, essentially $L(2\lambda,\chi^2)/L(2\lambda+1,\chi^2)$
(Lemma~\ref{lem:intertwining}). For $SL(2)$ it is essentially
$L(\lambda,\chi)/L(\lambda+1,\chi)$. The traces agree and the weights do
not. Under $(\chi,\lambda)\mapsto(\chi^2,2\lambda)$ the weights agree
exactly. This is the first failure found in~\cite{JC26}, and
Lemma~\ref{lem:intertwining} supplies the proof of the formula it
quotes.

\paragraph{The residual spectrum.} For the cover it is the sum of the
$V_b\cong r^+_{\psi_{-b}}$, $b\in k^\times/k^{\times2}$
(Theorem~\ref{thm:residual}), with exponent $\frac12$. For $SL(2)$ it is
the trivial representation, with exponent $1$. No $V_b$ corresponds
directly to a representation in the discrete spectrum of $SL(2)$. Under
the squaring correspondence every $V_b$ goes to the trivial
representation.

\paragraph{The term at $\lambda=0$.} On the spectral side, each character of the torus that is fixed by the
Weyl group contributes $\frac14\operatorname{tr}(M(w,0)I(\cdot,f))$. For
the cover this is Theorem~\ref{thm:Jchi}(c), and for $SL(2)$ it is the
same computation~\cite{LL79,GJ79}. The fixed characters correspond under the
direct comparison: they are $\chi_{\psi_b}=\chi_\psi\cdot(\kappa_b\circ p)$
for the cover (Lemma~\ref{lem:w-action}(d)) and the quadratic characters
$\kappa_b$ for $SL(2)$. The operators do not correspond. For the cover
$M(w,0)=-1$ for every $b$ (Proposition~\ref{prop:Mw0}), and the term is
$-\frac14\operatorname{tr}I(\chi_{\psi_b,0},f)$. For $SL(2)$, $M(w,0)=-1$
when $b=1$. When $b\ne1$ the representation $I^\flat(\kappa_b)$ is
reducible at every place, $I^\flat(\kappa_{b,v})=\pi^+_v\oplus\pi^-_v$, and
$M(w,0)$ acts on $\bigotimes_v\pi_v^{\varepsilon_v}$ by
$\prod_v\varepsilon_v$, for a suitable labeling~\cite{LL79}. So for
directly matching functions the two terms differ by
\begin{equation}\label{eq:lambda0}
\frac14\Bigl(\operatorname{tr}\bigl(M(w,0)I^\flat(\kappa_b,\phi)\bigr)+
\operatorname{tr}I^\flat(\kappa_b,\phi)\Bigr)=\frac12\sum_{\prod_v
\varepsilon_v=1}\ \prod_v\operatorname{tr}\pi_v^{\varepsilon_v}(\phi_v),
\end{equation}
for each $b\ne1$: half of the trace of $\phi$ on the part of
$I^\flat(\kappa_b)$ on which $M(w,0)=1$. In~\cite{LL79} this term is
accounted for by the torus of $k(\sqrt b)$: the $\kappa$-orbital
integrals lead to the trace formula of that torus, and~\eqref{eq:lambda0}
is the contribution of its trivial character. On the cover
nothing is left over at $\lambda=0$. The right side of~\eqref{eq:lambda0}
has the same shape as the trace on the theta functions,
$\frac12\Theta^{\mathrm{st}}_b(f)+\frac12\Theta^{\mathrm{st}}_b(f_-)=\sum_{
\prod\varepsilon_v=1}\prod_v\operatorname{tr}r^{\varepsilon_v}_{\psi'_v}(f_v)$
(Proposition~\ref{prop:theta-packet}(b)). On the cover the representations
with this shape are in the discrete spectrum, and on $SL(2)$ they are in
a term of the continuous spectrum.

\paragraph{Why multiplicities cannot be compared.} For $SL(2)$ the
multiplicity of a representation inside its $L$-packet is decided by the
unstable part of the trace formula. That part has two sources: the
$\kappa$-orbital integrals on the elliptic tori, and the operators
$M(w,0)$ at the quadratic characters. Both lead to the tori $T_E$ as
endoscopic groups~\cite{LL79}. The cover has neither. The
$\kappa$-orbital integral of $f$ is a stable orbital integral of $f_-$
(Lemma~\ref{lem:kappa-orbital}(c)), $M(w,0)$ is the scalar $-1$, and
both endoscopic groups are $SO(3)$. What decides multiplicities on the
cover is the pair $(f,f_-)$, that is, the central character and the
sign $\varepsilon(\frac12,\sigma)$ in Waldspurger's theorem. So the
information that would be needed to prove multiplicity one for $SL(2)$
sits in terms of its trace formula that have no image on the cover.
The text~\cite{JC26} reaches the same conclusion by a heuristic
argument: a comparison through $PGL(2)$ controls sums over $L$-packets
and not the multiplicity of one representation. The description above
says which terms are responsible.

\paragraph{The action of $GL(2)$ and the fibers of the squaring map.}
The squaring correspondence is two-to-one on unramified parameters:
$t'$ and $-t'$ have the same square. The text~\cite{JC26} also finds
that its bijection between characters of the torus and genuine
characters depends on a choice, with exactly two possibilities that
differ by the unramified quadratic character, neither of them
preferred. Both facts are the action of $GL(2)$ on the cover.

\begin{proposition}\label{prop:fibers}
\begin{enumerate}[label=(\alph*)]
\item (Local.) Let $\chi$ be a character of $F^\times$, $a\in F^\times$,
  and $\kappa_a=(a,\cdot)_v$. Then $I(\chi_{\psi}\cdot(\chi\circ p))\circ
  \Ad(\delta_a)\cong I(\chi_\psi\cdot(\chi\kappa_a\circ p))=I(
  \chi_{\psi_a}\cdot(\chi\circ p))$. If $v$ is odd, $\chi$ is unramified
  and $u$ is a unit that is not a square, then $\Ad(\delta_u)$ preserves
  $K^*$ and exchanges the unramified genuine representations with
  parameters $\chi(\varpi)$ and $-\chi(\varpi)$.
\item (Global.) The orbits of $GL(2,k)$ on the classes $(A,\mu)$,
  $\mu=\chi_\psi\cdot(\chi\circ p)$, are the sets on which $\chi^2$ is
  constant. In particular the residual representations $V_b$ form one
  orbit.
\end{enumerate}
\end{proposition}

\begin{proof}
(a) The automorphism $\Ad(\delta_a)$ preserves $\overline B$, $N$ and
$\overline A$, and it acts on $\overline A$ by multiplication by
$\kappa_a$ (Lemma~\ref{lem:Ad}(d)). So composing with it carries the
representation induced from $\mu$ to the one induced from $\mu\kappa_a$;
on characters this is Proposition~\ref{prop:IA-gl2}(c). The second form
is Lemma~\ref{lem:weil}(d). For the last statement use
Lemma~\ref{lem:Ad}(e) and $\kappa_u(\varpi)=-1$.

(b) By Lemma~\ref{lem:spectrum-gl2}(b), $h\in GL(2,k)$ sends $(A,\mu)$ to
$(A,\mu\kappa_a)$, $a=\det h$. Two characters $\chi,\chi'$ of
$k^\times\backslash\A_k^1$ have the same square exactly when
$\chi'=\chi\kappa_b$ for some $b\in k^\times$ (proof of
Lemma~\ref{lem:w-action}(d)), and $U_hV_b=V_{ab}$.
\end{proof}

So the fibers of the squaring map are the orbits of $GL(2)$. The direct
comparison needs a section of the squaring map, and there is none that
$GL(2)$ respects. Replacing $\psi$ by $\psi_u$ replaces $\Delta_\psi$ by
$\eta_x(u)\Delta_\psi$ (Lemma~\ref{lem:Delta}(d)) and exchanges the two
bijections of~\cite{JC26}. This is why neither is preferred. For the
cuspidal spectrum the corresponding statement is part of Waldspurger's
theorem: composing a genuine cuspidal representation with
$\Ad(\delta_a)$ twists its image in $PGL(2)$ by $\kappa_a$, when that
image is defined. We quote this and do not prove it
(Remark~\ref{rem:compare-status}).

\begin{example}\label{ex:fibers}
Let $k=\Q$ and $S=\{\infty,2\}$. The four residual representations that
are unramified outside $S$ are $V_b$, $b\in\{\pm1,\pm2\}$
(Example~\ref{ex:ramified-Q}). The matrices $\diag(b,1)$ permute them
transitively, and all four correspond to the trivial representation of
$SL(2,\A)$ under the squaring map, since $\kappa_b^2=1$. Under the direct
comparison they would correspond to four different representations
induced from the quadratic characters $\kappa_b$ at exponent
$\frac12$, none of which is in the discrete spectrum of $SL(2,\A)$.
\end{example}

\section{Status of the results}\label{sec:status-compare}

This section lists what in Chapter~\ref{sec:compare} is proved, what is
quoted, and what is not verified.

\begin{remarkrev}\label{rem:compare-status}
\emph{Proved here}, given the earlier chapters:
Lemmas~\ref{lem:direct-torus} and~\ref{lem:Delta-product},
Propositions~\ref{prop:wrong-orbit}, \ref{prop:unit-swap},
\ref{prop:epsilon} and~\ref{prop:fibers},
Corollaries~\ref{cor:coherent} and~\ref{cor:no-fl}, and the
relations~\eqref{eq:direct-char}--\eqref{eq:direct-ps}.
Corollary~\ref{cor:coherent} is conditional on the existence of matching
functions at every place, and Corollary~\ref{cor:no-fl} shows that this
condition cannot be met with the unit elements. The chapter was written
in one pass by Claude in November 2026 and has not been refereed.

\emph{Quoted and not rechecked}: the statements about the trace formula
of $SL(2)$ in Section~\ref{sec:compare-spectral}, namely the normalizing
factor $L(\lambda,\chi)/L(\lambda+1,\chi)$, the action of $M(w,0)$ on
$I^\flat(\kappa_b)$ and the role of~\eqref{eq:lambda0} in~\cite{LL79}.
They were written from memory of~\cite{LL79} and~\cite{GJ79}. The
statement about cuspidal representations after
Proposition~\ref{prop:fibers} is quoted from memory
of~\cite{Wa80,Wa91}. The identity~\eqref{eq:packet-lift}, used in
Section~\ref{sec:wrong-orbit}, has the status given in
Chapter~\ref{sec:stable}. All statements about the content
of~\cite{JC26} refer to its Draft~7.

\emph{Checked numerically} (script \texttt{comparison\_n2\_check.py}):
Lemma~\ref{lem:direct-torus} at $v=\infty,2,3,5,7$;
Lemma~\ref{lem:Delta-product} for random elements of $SL(2,\Q)$;
Proposition~\ref{prop:unit-swap}(a),~(d) and the fact that
$\Ad(\delta_u)$ preserves $K^*$, for $p=3,5,7,11,13$; and the values in
Examples~\ref{ex:coherent} and~\ref{ex:epsilon}. The orbital integrals in
Proposition~\ref{prop:unit-swap}(b),~(c) are not computed by the script.
\end{remarkrev}

\begin{example}\label{ex:status-compare}
The value $\varepsilon^S(\gamma)=-i$ of Example~\ref{ex:epsilon} is of the
third kind for its ingredients, the three values of $\Delta_{\psi_v}$, and
of the first kind for the conclusion, which is
Proposition~\ref{prop:epsilon}. The statement in the same example
about the two classes uses only that proposition and
Proposition~\ref{prop:unit-swap}(d).
\end{example}
